\chapter{Présentation des résultats}
\label{chap:resultats}

Ce chapitre présente les résultats de la campagne de simulation conduite sur la fonction
\texttt{computeRFM} d'ELMES-TEKA. La présentation suit le cadre \emph{ADEMP} proposé par
\citet{Morris2019} — objectifs, mécanismes générateurs, cibles, méthodes et mesures de
performance — dont les éléments constitutifs sont décrits à l'Annexe~\ref{ann:protocole}.
Le chapitre se limite à l'exposition factuelle des résultats ; leur interprétation est réservée
au chapitre suivant.

% ------------------------------------------------------------------
\section{Cibles mesurées et indicateurs retenus}
\label{sec:resultats-cibles}

Deux objectifs distincts structurent la campagne.

\paragraph{Conformité à la référence indépendante (H1).}
La cible principale est l'accord entre les sorties du moteur JavaScript \texttt{computeRFM}
et celles d'une implémentation de référence en Python, calculant les sommes monétaires en
centimes entiers. L'unité d'analyse est l'essai ; un essai est dit non-conforme dès qu'au
moins un client présente un écart, que ce soit sur un score ($s_R$, $s_F$ ou $s_M$) ou sur
l'étiquette de segment. Le taux de non-conformité est le rapport entre le nombre d'essais
non conformes et le nombre total d'essais.

\paragraph{Robustesse aux perturbations (H2).}
La cible secondaire est le taux de changement de segment parmi les clients communs à deux
partitions appariées — la partition initiale d'un corpus et la partition obtenue après
application d'une perturbation. Pour H2, la cellule d'intérêt primaire est définie par le
bruit monétaire $\pm1\,\%$ appliqué aux corpus de taille 50 ou 200, avec profil régulier ou
dissymétrique, en ne retenant que les clients éloignés des six seuils initiaux. Le critère
de compatibilité est la borne supérieure de l'intervalle de confiance Monte-Carlo à $95\,\%$
strictement inférieure ou égale à $5\,\%$.

\paragraph{Indicateurs calculés.}
Le dictionnaire complet des variables est fourni au tableau~\ref{tab:ann-dictionnaire}
(Annexe~\ref{ann:protocole}). Les indicateurs retenus pour la synthèse sont :

\begin{itemize}
  \item \texttt{change\_rate} — proportion de clients communs changeant de segment ;
  \item \texttt{far\_change\_rate} — même proportion, restreinte aux clients éloignés des
    seuils initiaux selon la règle à $5\,\%$ ;
  \item \texttt{lost\_rate} — proportion de clients présents dans le corpus initial mais
    absents après perturbation ;
  \item \texttt{oracle\_pass} — indicateur binaire d'absence de tout écart à la référence ;
  \item \texttt{occupied\_segments} — nombre de catégories non vides après traitement ;
  \item \texttt{ari} — indice de Rand ajusté entre les deux partitions des clients communs.
\end{itemize}

Les intervalles Monte-Carlo à $95\,\%$ sont calculés par loi $t$ à partir des 100 répétitions
par cellule, puis tronqués à $[0;1]$. Un intervalle de largeur nulle traduit une absence de
variation observée dans les répétitions de la cellule considérée ; il ne constitue pas une
démonstration d'une probabilité véritable de changement exactement nulle.

% ------------------------------------------------------------------
\section{Résultats de conformité à la référence indépendante}
\label{sec:resultats-conformite}

\subsection{Cas manuels}
\label{subsec:resultats-cas-manuels}

Neuf cas manuels ont été enregistrés dans \texttt{conformity.json} avant l'exécution de
la campagne. Ces cas couvrent les configurations limites suivantes : corpus vide, six clients
avec cinq groupes distincts aux frontières de quantiles, client singleton, valeurs toutes
identiques, filtrages par statut de commande (\texttt{CONFIRMED}, \texttt{PAID},
\texttt{DELIVERED}), exclusion par tenant et exercice, client absent du répertoire des
clients mais dont l'identifiant figure dans les commandes, montant nul, et traitement de la
partie entière d'une fraction de journée.

Les neuf cas sont satisfaits sans exception (tableau~\ref{tab:ann-cas-manuels},
Annexe~\ref{ann:protocole}). Ce résultat ne supprime pas les écarts observés dans d'autres
essais de la campagne principale.

\subsection{Campagne complète et premier contre-exemple}
\label{subsec:resultats-campagne}

La campagne porte sur 12\,000 essais produits par le croisement de quatre tailles de corpus
($n \in \{3, 10, 50, 200\}$ clients), trois profils de génération (\texttt{regular},
\texttt{skewed}, \texttt{ties}), 100 répétitions par cellule et 10 scénarios de perturbation.
Chaque perturbation est appariée au corpus initial correspondant ; les seuils de segmentation
sont recalculés après chaque perturbation.

\begin{table}[htb]
  \centering
  \caption{Bilan de conformité sur l'ensemble de la campagne.}
  \label{tab:resultats-conformite-bilan}
  \begin{tabular}{lr}
    \toprule
    Indicateur & Valeur \\
    \midrule
    Essais totaux                                       & 12\,000 \\
    Essais avec au moins un écart (référence centimes)  & 200 \\
    Taux de non-conformité                              & $1{,}667\,\%$ \\
    Occurrences de score $M$ erroné                     & 489 \\
    Occurrences d'étiquette de segment erronée          & 2 \\
    Exceptions moteur dans le domaine principal         & 0 \\
    Cas manuels satisfaits                              & 9 / 9 \\
    \bottomrule
  \end{tabular}
\end{table}

Sur les 200 essais non conformes, tous les écarts portent sur la dimension monétaire.
Les 489 occurrences de score $M$ erroné concernent des clients pour lesquels la somme des
montants de commandes diffère selon le mode de calcul : addition en virgule flottante dans
le moteur JavaScript contre addition exacte en centimes entiers dans la référence Python.
Les deux occurrences d'étiquette de segment erronée découlent de ces écarts de score $M$.
Aucun essai ne présente d'exception non gérée dans les scénarios appartenant au domaine
principal de la campagne.

Il est à noter que la référence de comparaison a été précisée en deux étapes : la première
implémentation en virgule flottante relevait 189 essais avec écart ; la reformulation en
centimes entiers en relève 200. Les diagnostics des deux étapes sont conservés dans l'archive.

\paragraph{Premier contre-exemple.}
Le premier essai non conforme enregistré est identifié par l'essai
\texttt{n3\allowbreak-ties\allowbreak-r16\allowbreak-amount\_noise\_1pct}, avec la graine de
corpus $20461118$ et la graine de perturbation $30461137$.

\begin{table}[htb]
  \centering
  \caption{Premier écart de score conservé (montants en USD).}
  \label{tab:resultats-contre-exemple}
  \begin{tabular}{llrcc}
    \toprule
    Client & Tickets (USD) & Somme (USD) & $M$ attendu & $M$ obtenu \\
    \midrule
    00000 & $9{,}99$ ; $10{,}00$ ; $9{,}97$ & $29{,}96$ & 1 & 1 \\
    00001 & $10{,}03$ ; $9{,}98$ ; $10{,}04$ & $30{,}05$ & 1 & 1 \\
    00002 & $9{,}96$ ; $10{,}02$ ; $10{,}07$ & $30{,}05$ & 1 & 2 \\
    \bottomrule
  \end{tabular}
\end{table}

Les clients 00001 et 00002 présentent des sommes affichées identiques ($30{,}05$~USD), mais
la représentation interne en virgule flottante produit des résultats légèrement différents
après addition successive des trois tickets. La référence en centimes considère les deux
sommes comme égales et attribue à chaque client le score $M = 1$ ; le moteur JavaScript
en obtient des valeurs distinctes et attribue le score $M = 2$ au client 00002. Les trois
étiquettes de segment demeurent « À risque » dans cet essai.

% ------------------------------------------------------------------
\section{Résultats des mesures de performance par scénario}
\label{sec:resultats-performance}

\subsection{Critère principal de robustesse H2}
\label{subsec:resultats-h2}

Le critère H2 porte sur les cellules vérifiant les trois conditions suivantes : taille de
corpus $n \in \{50, 200\}$, profil régulier ou dissymétrique, et scénario de bruit monétaire
$\pm1\,\%$. Pour chaque client commun, six distances normalisées aux seuils initiaux sont
calculées ; un client est dit éloigné si chacune de ces six distances est supérieure à
$5\,\%$. Seules les répétitions conservant au moins dix clients éloignés communs aux deux
partitions sont admissibles pour H2.

\begin{table}[htb]
  \centering
  \caption{Résultats du critère H2 — bruit monétaire $\pm1\,\%$, clients éloignés des seuils.}
  \label{tab:resultats-h2}
  \begin{tabular}{rlrrrl}
    \toprule
    Taille & Profil & Rép. admissibles & Taux moyen (\%) & IC $95\,\%$ (\%) & Décision \\
    \midrule
    50  & \texttt{regular} & 100 & $0{,}000$ & $[0{,}000;\,0{,}000]$ & Compatible \\
    50  & \texttt{skewed}  & 100 & $0{,}000$ & $[0{,}000;\,0{,}000]$ & Compatible \\
    200 & \texttt{regular} & 100 & $0{,}000$ & $[0{,}000;\,0{,}000]$ & Compatible \\
    200 & \texttt{skewed}  & 100 & $0{,}000$ & $[0{,}000;\,0{,}000]$ & Compatible \\
    \bottomrule
  \end{tabular}
\end{table}

Dans les quatre cellules admissibles pour H2, le taux de changement observé parmi les clients
éloignés est nul dans la totalité des répétitions, ce qui produit des intervalles de confiance
ponctuels $[0;0]$. La borne supérieure de ces intervalles est strictement inférieure au seuil
de $5\,\%$ ; le critère H2 est déclaré compatible dans chacune de ces cellules. Les détails
complets par répétition sont disponibles dans le fichier \texttt{hypothesis\_H2.csv} de
l'archive expérimentale.

\subsection{Taux de changement — synthèse par taille et profil}
\label{subsec:resultats-synthese}

Le tableau~\ref{tab:resultats-synthese} présente les taux moyens de changement de segment
pour les profils \texttt{regular} et \texttt{skewed} à la taille $n = 200$, ainsi que pour
le profil \texttt{ties} à la même taille. Les données complètes pour toutes les combinaisons
taille/profil (120 lignes) figurent à l'Annexe~\ref{ann:reproduction}.

\begin{table}[htb]
  \centering
  \caption{Synthèse des taux de changement, taille $n = 200$. Moyenne sur 100 répétitions.
    IC : intervalle Monte-Carlo à $95\,\%$. Profils \texttt{regular} et \texttt{skewed}
    présentés en plage ; profil \texttt{ties} séparé.}
  \label{tab:resultats-synthese}
  \small
  \begin{tabular}{lcc}
    \toprule
    Scénario
      & \texttt{regular} / \texttt{skewed} — changement (\%)
      & \texttt{ties} — changement (\%) \\
    \midrule
    Initial
      & $0{,}000$
      & $0{,}000$ \\
    Montants $\pm1\,\%$
      & $0{,}035$ à $0{,}055$
      & $0{,}000$ \\
    Montants $\pm5\,\%$
      & $0{,}150$ à $0{,}265$
      & $0{,}000$ \\
    Montants $\pm20\,\%$
      & $0{,}640$ à $0{,}925$
      & $0{,}000$ \\
    Dates $\pm1$\,j
      & $0{,}970$ à $1{,}035$
      & $49{,}760$ \\
    Dates $\pm7$\,j
      & $4{,}660$ à $4{,}685$
      & $63{,}375$ \\
    Doublons $1\,\%$
      & $1{,}000$ à $1{,}275$
      & $0{,}000$ \\
    Doublons $5\,\%$
      & $3{,}655$ à $4{,}120$
      & $0{,}000$ \\
    Manquantes $1\,\%$
      & $1{,}896$ à $1{,}921$
      & $0{,}000$ \\
    Manquantes $5\,\%$
      & $6{,}661$ à $7{,}395$
      & $0{,}000$ \\
    \bottomrule
  \end{tabular}
\end{table}

Pour les profils \texttt{regular} et \texttt{skewed}, les taux de changement augmentent avec
l'intensité de la perturbation. Pour le profil \texttt{ties}, les taux demeurent nuls pour
toutes les perturbations monétaires et nuls ou quasi-nuls pour les perturbations de doublons
et de commandes manquantes, tandis qu'ils atteignent des valeurs élevées pour les perturbations
de dates (voir section~\ref{subsec:resultats-ties}).

La figure~\ref{fig:ann-sensibilite-montants} (Annexe~\ref{ann:reproduction}) présente les
moyennes et intervalles Monte-Carlo à $95\,\%$ pour les perturbations monétaires, sur
l'ensemble des quatre tailles et des trois profils. La figure~\ref{fig:ann-transitions}
(Annexe~\ref{ann:reproduction}) illustre les transitions individuelles de segment pour
$n = 200$, profil \texttt{regular}, bruit monétaire $\pm5\,\%$, agrégées sur 100 répétitions.

\subsection{Taux de clients perdus}
\label{subsec:resultats-perdus}

Le taux de clients perdus (\texttt{lost\_rate}) mesure la proportion de clients présents
dans le corpus initial mais absents après perturbation. Ce taux est nul pour toutes les
perturbations monétaires et de dates, indépendamment de la taille et du profil. Des clients
perdus apparaissent uniquement dans les scénarios de suppression de commandes (\emph{manquantes}).

\begin{table}[htb]
  \centering
  \caption{Taux de clients perdus pour les scénarios de suppression de commandes,
    sélection de cellules représentatives.}
  \label{tab:resultats-perdus}
  \begin{tabular}{rlcc}
    \toprule
    Taille & Profil & Manquantes $1\,\%$ (\%) & Manquantes $5\,\%$ (\%) \\
    \midrule
    3   & \texttt{regular}  & $0{,}000$ & $1{,}333$ \\
    3   & \texttt{skewed}   & $0{,}333$ & $0{,}333$ \\
    10  & \texttt{regular}  & $0{,}000$ & $0{,}200$ \\
    10  & \texttt{skewed}   & $0{,}000$ & $0{,}400$ \\
    50  & \texttt{regular}  & $0{,}100$ & $0{,}440$ \\
    50  & \texttt{skewed}   & $0{,}100$ & $0{,}620$ \\
    200 & \texttt{regular}  & $0{,}070$ & $0{,}475$ \\
    200 & \texttt{skewed}   & $0{,}080$ & $0{,}495$ \\
    \bottomrule
  \end{tabular}
\end{table}

Pour le profil \texttt{ties}, le taux de clients perdus est nul dans les scénarios de
suppression à $1\,\%$, et inférieur à $0{,}05\,\%$ pour la suppression à $5\,\%$ (tailles
50 et 200). Les valeurs complètes figurent à l'Annexe~\ref{ann:reproduction}.

\subsection{Comportement du profil identique (\emph{ex~æquo})}
\label{subsec:resultats-ties}

Le profil \texttt{ties} est construit de telle sorte que tous les clients présentent
exactement le même nombre de commandes (trois), la même date de dernier achat (30~jours
avant l'observation) et le même ticket (10~USD). Dans la situation initiale, une seule
catégorie de segment est occupée.

Ce profil produit deux comportements distincts selon le type de perturbation :

\begin{enumerate}
  \item \textbf{Perturbations monétaires et de doublons} : les taux de changement sont nuls
    dans toutes les répétitions et pour toutes les tailles testées. La perturbation monétaire
    ne modifie pas les relations d'ordre entre clients à valeurs identiques ; l'ajout de
    commandes dupliquées ne rompt pas non plus les égalités de fréquence.

  \item \textbf{Perturbations de dates} : les taux de changement sont élevés et présentent
    des intervalles de confiance étroits du fait de leur faible variabilité entre répétitions.
    Les valeurs observées pour les quatre tailles sont résumées dans le
    tableau~\ref{tab:resultats-ties-dates}.
\end{enumerate}

\begin{table}[htb]
  \centering
  \caption{Taux de changement du profil \texttt{ties}, perturbations de dates.}
  \label{tab:resultats-ties-dates}
  \begin{tabular}{rcccc}
    \toprule
    & \multicolumn{2}{c}{Dates $\pm1$\,j} & \multicolumn{2}{c}{Dates $\pm7$\,j} \\
    \cmidrule(lr){2-3}\cmidrule(lr){4-5}
    $n$ & Moyen (\%) & IC $95\,\%$ (\%) & Moyen (\%) & IC $95\,\%$ (\%) \\
    \midrule
    3   & $48{,}333$ & $[43{,}887;\,52{,}780]$ & $63{,}333$ & $[61{,}129;\,65{,}538]$ \\
    10  & $42{,}200$ & $[39{,}155;\,45{,}245]$ & $56{,}800$ & $[55{,}641;\,57{,}959]$ \\
    50  & $48{,}640$ & $[46{,}055;\,51{,}225]$ & $62{,}500$ & $[61{,}871;\,63{,}129]$ \\
    200 & $49{,}760$ & $[46{,}761;\,52{,}759]$ & $63{,}375$ & $[62{,}877;\,63{,}873]$ \\
    \bottomrule
  \end{tabular}
\end{table}

Le profil \texttt{ties} avec perturbation de dates constitue un cas de figure distinct des
profils \texttt{regular} et \texttt{skewed}. Les données complètes par scénario et taille
pour ce profil figurent à l'Annexe~\ref{ann:reproduction}.

% ------------------------------------------------------------------
\section{Observations hors domaine et limites du banc}
\label{sec:resultats-hors-domaine}

Cinq configurations situées hors du périmètre de la campagne principale ont été soumises à
la fonction instrumentée. Ces configurations ne font pas partie des 12\,000 essais et ne
contribuent pas aux mesures de performance.

\begin{table}[htb]
  \centering
  \caption{Comportements observés dans cinq configurations hors domaine.}
  \label{tab:resultats-hors-domaine}
  \begin{tabular}{p{0.26\textwidth}p{0.32\textwidth}p{0.32\textwidth}}
    \toprule
    Entrée & Comportement observé & Portée de l'observation \\
    \midrule
    USD et CDF mélangés
      & Valeurs additionnées sans séparation : $10 + 28\,000 = 28\,010$
      & Absence de conversion ou de séparation de devise à ce niveau de la fonction \\
    Identifiant client absent
      & Exception \texttt{TypeError}
      & Rejet métier contrôlé non démontré \\
    Date invalide
      & Exception \texttt{RangeError}
      & Rejet métier contrôlé non démontré \\
    Date future ($6$ oct. 2026, observation au $1$\textsuperscript{er} oct.)
      & Acceptée sans erreur
      & Cas hors domaine ; comportement dans d'autres couches du SaaS non vérifié ici \\
    Montant négatif ($-10$)
      & Accepté sans erreur
      & Politique de remboursements non définie dans ce contrôle \\
    \bottomrule
  \end{tabular}
\end{table}

Les exceptions \texttt{TypeError} et \texttt{RangeError} signalent que la fonction réagit
à ces entrées, mais elles ne démontrent pas l'existence d'un contrôle métier explicite en
amont. Ces constats sont indépendants des vérifications effectuées dans les autres couches
du SaaS ELMES-TEKA.

\paragraph{Limites de la campagne.}
Les points suivants délimitent la portée des résultats présentés dans ce chapitre :

\begin{itemize}
  \item Les corpus sont synthétiques et non calibrés sur des données de petites entreprises
    congolaises réelles ; les résultats décrivent le comportement de la fonction sur ces
    mécanismes générateurs, et non son comportement attendu en production.
  \item Les intervalles Monte-Carlo de largeur nulle indiquent une absence de variation
    dans les répétitions observées ; ils ne constituent pas une preuve d'une probabilité
    de changement identiquement nulle.
  \item L'intégration MongoDB, les routes HTTP et les effets économiques de l'application
    n'ont pas été évalués dans cette campagne.
  \item Aucune vérification MongoDB réelle ni reproduction indépendante n'ont été effectuées
    pour préparer les présents résultats.
\end{itemize}

% ------------------------------------------------------------------
\section{Traçabilité des résultats}
\label{sec:resultats-tracabilite}

L'ensemble des résultats présentés dans ce chapitre est dérivé de l'archive expérimentale
\texttt{ELMES\_TEKA\_Resultats\_Mission\_RFM\_2026-11-02.zip}, dont l'intégrité a été
vérifiée par empreinte SHA-256 pour 55~fichiers de livraison et 12~fichiers de données
brutes. Les 12\,000~enregistrements JSON ont été relus et comptés.

\begin{table}[htb]
  \centering
  \caption{Environnement d'exécution consigné dans l'archive expérimentale.}
  \label{tab:resultats-versions}
  \begin{tabular}{ll}
    \toprule
    Composant & Version \\
    \midrule
    Node.js               & 24.22.1 \\
    TypeScript            & 5.9.3 \\
    Python                & 3.12.14 \\
    NumPy                 & 2.3.5 \\
    pandas                & 2.2.3 \\
    SciPy                 & 1.17.0 \\
    Matplotlib            & 3.10.8 \\
    Commit applicatif     & \texttt{\small b062a93a7be96d082dee9afe66b9f4f1f9ebecaa} \\
    Révision du banc      & \texttt{\small a4dbbafba0724a269482706b94288bb28cd5db3d} \\
    \bottomrule
  \end{tabular}
\end{table}

La procédure de reproduction à partir d'une copie complète du dépôt au commit pertinent
est documentée à l'Annexe~\ref{ann:reproduction}, section~«~Versions, intégrité et reproduction~». Le lien vers l'archive complète est soumis à l'accès au compte de stockage
de l'auteur ; une solution de conservation accessible au jury reste à convenir.

Le registre \texttt{PROVENANCE\_ANNEXES.md} (Annexe~\ref{ann:protocole}) précise la
correspondance entre chaque donnée présentée dans ce chapitre et le fichier source dans
l'archive expérimentale.
